\section{Infrastructure}
\label{appx:infra}
We cover additional details about how \bugs{} works, specifically
\begin{itemize}
    \item The form factor of a \bugs{} task instance.
    \item How we identify repositories and the SWE-agent configuration we use to automatically install them.
    \item How the task validation and evaluation harnesses work.
\end{itemize}

\subsection{\bugs{} Task Instance}
\label{appx:infra:anatomy}
We briefly review the format of a \bugs{} task instance, highlight how it is different from a SWE-bench task instance, and discuss why \bugs{}'s relatively simple infrastructure compared to SWE-bench allows us scale task collection much more efficiently.

A \bugs{} task instance is very similar to the form factor of a SWE-bench task instance, with several minor differences.
A \bugs{} task instance includes the following fields:
\begin{itemize}
    \item \texttt{repo}: The repository the task instance is from.
    \item \texttt{instance\_id}: A unique identifier (usually \texttt{(repo).(bug\_type).(hash)})
    \item \texttt{base\_commit}: Hash of the GitHub branch that points to the repository with the bug \texttt{patch} applied.
    \item \texttt{patch}: The \texttt{diff} that causes the bug. It is applied to the original codebase to create the bug. Reverting this patch is effectively the solution.
    \item \texttt{problem\_statement}: The generated issue text that conveys the bug.
    It is provided to a model or system before it begins attempting a fix.
    \item \texttt{created\_at}: A timestamp matching when the bug was successfully validated and pushed to the mirror repository as a branch.
    \item \texttt{FAIL\_TO\_PASS}: The unit tests that break when the test suite is run with the bug \texttt{patch} applied.
    \item \texttt{PASS\_TO\_PASS}: The unit tests that do not break.
    These correspond to the set of all tests minus the \texttt{FAIL\_TO\_PASS} tests.
\end{itemize}

We summarize the key distinctions between a \bugs{} and SWE-bench task instance:
\begin{itemize}
    \item \bugs{} task instances do not include the \texttt{version} or \texttt{environment\_setup\_commit} fields, which SWE-bench requires as additional identifiers for specifying repository-specific installation instructions across time.
    In \bugs{}, unique installation instructions are specified for each (repository, commit).
    \item The \texttt{hints\_text} field is not included. In SWE-bench, this refers to the issue and PR thread comments written after the first commit of the corresponding PR.
    \item The \texttt{created\_at} field is assigned the timestamp reflecting when the bug was successfully validated.
    Originally, \texttt{created\_at} refers to when a PR was created.
    \item There is no \texttt{test\_patch} field, as the \bugs{} collection pipeline does not create or synthesize any hidden tests.
    All \texttt{FAIL\_TO\_PASS} bugs are visible and runnable in the repository at inference time.
\end{itemize}

\subsection{Repository Selection}
\label{appx:infra:repo_selection}
In addition to the criteria discussed in Section~\ref{sec:swesmith:collection}, we also ensure that a repository has a license that allows non-proprietary use.
The majority of software licenses are permissive (BSD, MIT, Apache), while the remainder are largely protective licenses (GPL) that still allow for non-commercial use.
We inspected the repositories with custom licenses and confirmed they allowed for the use cases exercised in our work.
The licenses for each repository are fully listed in Table~\ref{tab:licenses}.

\begin{table}[ht!]
    \centering
    \resizebox{\textwidth}{!}{%
\newcolumntype{P}[1]{>{\raggedright\arraybackslash}p{#1}} % Define a new left-aligned p column
\begin{tabular}{P{0.15\textwidth} | P{0.85\textwidth}}

\toprule
Apache License 2.0 & \footnotesize\texttt{Project-MONAI/MONAI; alanjds/drf-nested-routers; arrow-py/arrow; buriy/python-readability; facebookresearch/fvcore; getmoto/moto; google/textfsm; iterative/dvc; jax-ml/jax; jd/tenacity; kayak/pypika; modin-project/modin; pyca/pyopenssl; spulec/freezegun; tkrajina/gpxpy; tornadoweb/tornado; weaveworks/grafanalib} \\
\midrule
BSD 2-Clause "Simplified" License & \footnotesize\texttt{madzak/python-json-logger; pyasn1/pyasn1; pygments/pygments; sunpy/sunpy} \\
\midrule
BSD 3-Clause "New" or "Revised" License & \footnotesize\texttt{Suor/funcy; alecthomas/voluptuous; andialbrecht/sqlparse; cookiecutter/cookiecutter; dask/dask; django/channels; django/daphne; encode/starlette; gawel/pyquery; gweis/isodate; john-kurkowski/tldextract; lepture/mistune; oauthlib/oauthlib; pallets/click; pallets/flask; pallets/jinja; pallets/markupsafe; pandas-dev/pandas; scrapy/scrapy; theskumar/python-dotenv} \\
\midrule
GNU General Public License v3.0 & \footnotesize\texttt{Cog-Creators/Red-DiscordBot; adrienverge/yamllint} \\
\midrule
GNU Lesser General Public License v2.1 & \footnotesize\texttt{chardet/chardet; paramiko/paramiko; pylint-dev/astroid} \\
\midrule
GNU Lesser General Public License v3.0 & \footnotesize\texttt{Knio/dominate} \\
\midrule
ISC License & \footnotesize\texttt{kennethreitz/records} \\
\midrule
MIT License & \footnotesize\texttt{amueller/word\_cloud; borntyping/python-colorlog; bottlepy/bottle; cantools/cantools; cdgriffith/Box; cknd/stackprinter; conan-io/conan; cool-RR/PySnooper; datamade/usaddress; dbader/schedule; erikrose/parsimonious; facebookresearch/hydra; facelessuser/soupsieve; getnikola/nikola; graphql-python/graphene; hukkin/tomli; jaraco/inflect; jawah/charset\_normalizer; joke2k/faker; keleshev/schema; life4/textdistance; luozhouyang/python-string-similarity; marshmallow-code/apispec; marshmallow-code/marshmallow; marshmallow-code/webargs; martinblech/xmltodict; matthewwithanm/python-markdownify; mewwts/addict; mido/mido; mozillazg/python-pinyin; msiemens/tinydb; pdfminer/pdfminer; pndurette/gTTS; pudo/dataset; pydantic/pydantic; pyparsing/pyparsing; pytest-dev/iniconfig; python-hyper/h11; python-jsonschema/jsonschema; python-openxml/python-docx; pyupio/safety; pyvista/pyvista; r1chardj0n3s/parse; rsalmei/alive-progress; rubik/radon; rustedpy/result; scanny/python-pptx; seatgeek/thefuzz; sloria/environs; sqlfluff/sqlfluff; termcolor/termcolor; tobymao/sqlglot; tox-dev/pipdeptree; tweepy/tweepy; un33k/python-slugify; vi3k6i5/flashtext} \\
\midrule
Other & \footnotesize\texttt{Mimino666/langdetect; PyCQA/flake8; agronholm/exceptiongroup; agronholm/typeguard; aio-libs/async-timeout; benoitc/gunicorn; cloudpipe/cloudpickle; davidhalter/parso; django-money/django-money; gruns/furl; kurtmckee/feedparser; lincolnloop/python-qrcode; mahmoud/boltons; mahmoud/glom; mozilla/bleach; pexpect/ptyprocess; prettytable/prettytable; pwaller/pyfiglet; pydata/patsy; pydicom/pydicom; python-trio/trio; python/mypy; pyutils/line\_profiler; seperman/deepdiff} \\
\bottomrule
    \end{tabular}
    }
    \caption{License associated with each repository as of April 8, 2025. All licenses are permissive and allow for public, nonprofit use.}
    \label{tab:licenses}
\end{table}

We deliberately limit the search scope for repositories to those predominantly written in Python.
Following precedents, focusing on Python repositories allowed us to form assumptions about installation and testing procedures (e.g. repository is organized as a PyPI package, \texttt{pytest} is the testing framework) that made scaling up automatic repository setup with SWE-agent more tractable.
A worthwhile direction to consider for future work is expanding the coverage of repositories to be more comprehensive of codebases written in different programming languages, as ~\citet{yang_swe-bench_2024} does, extending SWE-bench style evaluation to JavaScript repositories with multimodal inputs.

\textbf{Automated repository installation.}
The goal of this step is to first, get the installation and testing instructions for a repository, and second, create a Docker image containing the repository with the development environment set up.

We provide the system prompt given to SWE-agent that asks it to install a repository in Figure~\ref{fig:prompt_install_repo}.
Each repository installation task is initialized with a clone of the original repository.
No additional steps (e.g. \texttt{pypi} package downloads, \texttt{conda} environment setup) are performed.

We run SWE-agent with \texttt{claude-3-5-sonnet-20241022} with a maximum cost limit of \$$2$ and a maximum call limit of $150$.
The installation run terminates whenever one of these conditions is met.
For every run, we record the interactions.
We then manually review the trajectory, identifying the appropriate installation and testing specifications.

Each run incurs an average cost of \$$0.72$ and an average of $17$ steps before SWE-agent issues the \texttt{submit} command.
The runs typically finish within two minutes.
The majority of Python repositories require fewer steps --- typically, SWE-agent will view the \texttt{CONTRIBUTING.md}, run the installation command provided verbatim in the text, and then runs \texttt{pytest}, showing all tests passing.
A minority of repositories will require several steps because additional dependencies must be installed with \texttt{apt-get}.
The manual review process following this requires $3$ to $20$ minutes.
One author carried out this effort for $128$ repositories, taking an estimated $18$ human hours to accomplish.
In the process of reaching $128$ repositories, the author gave up on $17$ repositories at the manual review stage.

\begin{example}[
System prompt for generating bugs with an LM
]
\small
\texttt{$<$uploaded\_files$>$} \\
\texttt{\{\{working\_dir\}\}} \\
\texttt{$<$/uploaded\_files$>$} \\
I've uploaded a python code repository in the directory \texttt{\{\{working\_dir\}\}}. \\

Can you please install this repository?
Your goal should be to configure the repository's development environment such that existing tests pass.
You are currently in the root directory of the repository, and nothing has been installed yet.
You in an Ubuntu 22.04 environment. \\

The repository is predominantly written in Python. Here are several tips for installing it: \\

1. A good place to start is to look for a \texttt{CONTRIBUTING.[md$\vert$rst]} file, which will often contain instructions on how to install the repository and any dependencies it may have. Occasionally, the \texttt{README.md} file may also contain installation instructions. \\

2. Usually, a repository may have \texttt{setup.py} or \texttt{pyproject.toml} files which can be used to install the package. \texttt{pip install -e .} is commonly used, although many packages will also require an additional specifier that installs development packages as well (e.g. \texttt{pip install -e .[dev]}). \\

3. To check whether the repository was installed successfully, run tests and see if they pass. You can usually find tests in a \texttt{tests/} or \texttt{test/} directory. You can run tests using \texttt{pytest} or \texttt{unittest}, depending on the framework used by the repository. \\

4. Sometimes, you will need to install additional packages, often listed in a \texttt{requirements.txt} or \texttt{environment.yml} file. Also, be mindful of Ubuntu system dependencies that may need to be installed via \texttt{apt-get} (e.g. \texttt{sudo apt-get install $<$package$>$}). \\

Once you are finished with installing the repository, run the \texttt{submit} command to submit your changes for review.
\end{example}
\captionof{figure}{
Prompt provided to SWE-agent + an LM asking it to install a repository.
}
\label{fig:prompt_install_repo}

\subsection{Validation, Evaluation Harnesses}
\label{appx:infra:harnesses}
We adapt SWE-bench's validation script to convert each bug patch into a SWE-bench style task instance.
This step ensures \bugs{} can be run by existing SWE-bench solutions.
The conversion involves two steps.
First, the bug patch is applied and pushed as a branch to a mirror clone of the repository.
Second, we create a SWE-bench style task instance from the bug patch, populating important fields such as Fail-to-Pass and Pass-to-Pass tests with information from the validation logs.

